% §1 Introduction
Evaluating a conversational agent outside production means talking to it with a simulator.
LLM-based user simulators now drive agent evaluation across task-oriented dialogue, tool use,
and agentic benchmarks \citep{sim2realAgentic2026,realUserSim2026,sierra2024taubench}.
The evaluation question every team asks is no longer ``did the agent pass'' but
``what kinds of conversations did we never even run'' ---
and no published method answers it measurably.

The raw material for an answer already exists on every deployment.
Production agents emit traces in standardized observability formats: the OpenTelemetry GenAI
semantic conventions (still marked Development as of October 2026 \citep{otelgenai2026}), the
OpenInference attribute specification \citep{openinference2026}, and LangSmith run trees
\citep{langsmith2026}.
Tooling mines these traces for failure prediction and benchmark construction
\citep{automataTraces2026,tracedance2026}, visualizes them as design-time graphs
\citep{graphoftrace2026}, and projects chat corpora into corpus-level maps
\citep{wildvis2024}.
But the conventions deliberately stop at the operation level: \texttt{gen\_ai.conversation.id}
is a linking key, not a conversation object.
The conversation-scale semantics --- which states the dialogue moved through, how it ended,
how deeply it regressed --- are left undefined by everyone.

Simulator validation, meanwhile, is measured at the wrong altitude for this question.
The 2026 sim2real wave compares response distributions in text space:
the distributional gap of \citet{distributionalGap2026} is response-level and text-only;
\citet{cds2026} scores distributional agreement of conversation-level statistics;
\citet{cue2026} calibrates user-simulator embeddings;
\citet{mirrorbench2026} benchmarks human-likeness of user proxies.
All of these ask ``do simulated responses look like real ones''.
None asks the structural question a test suite needs: \emph{which regions of conversational
behavior does the simulator exercise, and which real regions does it never touch?}
A simulator can match response distributions while never once producing the
clarify-then-abandon structure that a production cohort exhibits on bad-network days.

This paper joins the two artifacts.
We define a formal extraction mapping $\tau$ from standard observability exports to grammar
walks over the trajectory grammar $G=(Q,\Sigma,\delta)$ of our companion work
\citep{garcia2026grammar}, with terminal outcome classes and cycle depth as drift statistic.
We then define \emph{region coverage}: project real and simulated walk collections into a
2-D behavior plane whose basis is fit on the real set only, estimate where real conversations
live with a deterministic density grid, and measure which occupied regions the simulated set
reaches --- overall and per outcome class --- while the top-$K$ uncovered regions return their
nearest real exemplars: \emph{the tests you are not running}, made concrete and actionable.
The two clouds in one plane constitute the \emph{trajectory atlas}, the visualization and
metric layer for sim-versus-real transfer at the trajectory level.

Our contributions are:
\begin{enumerate}
  \item a formal extraction mapping $\tau$ from standard agent-observability formats
        (OTel GenAI, OpenInference, LangSmith) to trajectory-grammar walks with terminal
        outcome classes and cycle depth, shipped as a rule-based baseline whose classifier
        configuration is fully exposed, with an emission-model assumption stated so any
        calibrated state estimator can replace it without changing the interface
        (Section~\ref{sec:extraction});
  \item region-coverage evaluation: deterministic 2-D projection anchored on production,
        density-occupied real regions under a shared relative threshold, coverage
        overall and per outcome class, and top-$K$ uncovered regions with nearest real
        exemplars (Section~\ref{sec:coverage});
  \item the trajectory atlas as the visualization and metric layer for sim-versus-real
        comparison at the trajectory level, with a worked example on shipped fixtures and an
        evaluation protocol --- sensitivity sweep, class-conditional coverage, and a baseline
        ladder --- for paper-grade results (Sections~\ref{sec:atlas}
        and~\ref{sec:evaluation}).
\end{enumerate}

The paper is organized as follows.
Section~\ref{sec:background} recaps the grammar formalism, the distributional evaluation
lineage, and the verified facts about observability formats.
Section~\ref{sec:extraction} defines $\tau$.
Section~\ref{sec:coverage} defines region coverage.
Section~\ref{sec:atlas} presents the atlas and the worked example.
Section~\ref{sec:evaluation} gives the evaluation protocol.
Sections~\ref{sec:open_problems}--\ref{sec:conclusion} discuss open problems, limitations,
and conclusions.
